%\title{University of Bristol Thesis Template}
\RequirePackage[l2tabu]{nag}		% Warns for incorrect (obsolete) LaTeX usage
%
%
% File: memoirthesis.tex
% Author: Victor Brena
% Description: Contains the thesis template using memoir class,
% which is mainly based on book class but permits better control of 
% chapter styles for example. This template is an adaptation and 
% modification of Oscar's.
% 
% Memoir is a flexible class for typesetting poetry, fiction, 
% non-fiction and mathematical works as books, reports, articles or
% manuscripts. CTAN repository is found at:
% http://www.ctan.org/tex-archive/macros/latex/contrib/memoir/
%
%
% UoB guidelines for thesis presentation were found at:
% http://www.bris.ac.uk/esu/pg/pgrcop11-12topic.pdf#page=49
%
% UoB guidelines:
%
% The dissertation must be printed on A4 white paper. Paper up to A3 may be used
% for maps, plans, diagrams and illustrative material. Pages (apart from the
% preliminary pages) should normally be double-sided.
%
% Memoir class loads useful packages by default (see manual).
\documentclass[a4paper,11pt,leqno,openbib]{memoir} %add 'draft' to turn draft option on (see below)
\setlength{\parindent}{0em}
\setlength{\parskip}{1em}
%
%
% Adding metadata:
\usepackage{datetime}
\usepackage{ifpdf}
\ifpdf
\pdfinfo{
   /Author (Eddy Bertoluzzo)
   /Title (Master Thesis)
   /Keywords (One; Two;Three)
   /CreationDate (D:\pdfdate)
}
\fi
% When draft option is on. 
\ifdraftdoc 
	\usepackage{draftwatermark}				%Sets watermarks up.
	\SetWatermarkScale{0.3}
	\SetWatermarkText{\bf Draft: \today}
\fi
%
% Declare figure/table as a subfloat.
\newsubfloat{figure}
\newsubfloat{table}
% Better page layout for A4 paper, see memoir manual.
\settrimmedsize{297mm}{210mm}{*}
\setlength{\trimtop}{0pt} 
\setlength{\trimedge}{\stockwidth} 
\addtolength{\trimedge}{-\paperwidth} 
\settypeblocksize{634pt}{448.13pt}{*} 
\setulmargins{4cm}{*}{*} 
\setlrmargins{*}{*}{1.5} 
\setmarginnotes{17pt}{51pt}{\onelineskip} 
\setheadfoot{\onelineskip}{2\onelineskip} 
\setheaderspaces{*}{2\onelineskip}{*} 
\checkandfixthelayout
%
\frenchspacing
% Font with math support: New Century Schoolbook
\usepackage{fouriernc}
\usepackage[T1]{fontenc}
%
% UoB guidelines:
%
% Text should be in double or 1.5 line spacing, and font size should be
% chosen to ensure clarity and legibility for the main text and for any
% quotations and footnotes. Margins should allow for eventual hard binding.
%
% Note: This is automatically set by memoir class. Nevertheless \OnehalfSpacing 
% enables double spacing but leaves single spaced for captions for instance. 
\OnehalfSpacing 
%
% Sets numbering division level
\setsecnumdepth{subsection} 
\maxsecnumdepth{subsubsection}
%
% Chapter style (taken and slightly modified from Lars Madsen Memoir Chapter 
% Styles document
\usepackage{calc,soul,fourier}
\makeatletter 
\newlength\dlf@normtxtw 
\setlength\dlf@normtxtw{\textwidth} 
\newsavebox{\feline@chapter} 
\newcommand\feline@chapter@marker[1][4cm]{%
	\sbox\feline@chapter{% 
		\resizebox{!}{#1}{\fboxsep=1pt%
			\colorbox{gray}{\color{white}\thechapter}% 
		}}%
		\rotatebox{90}{% 
			\resizebox{%
				\heightof{\usebox{\feline@chapter}}+\depthof{\usebox{\feline@chapter}}}% 
			{!}{\scshape\so\@chapapp}}\quad%
		\raisebox{\depthof{\usebox{\feline@chapter}}}{\usebox{\feline@chapter}}%
} 
\newcommand\feline@chm[1][4cm]{%
	\sbox\feline@chapter{\feline@chapter@marker[#1]}% 
	\makebox[0pt][c]{% aka \rlap
		\makebox[1cm][r]{\usebox\feline@chapter}%
	}}
\makechapterstyle{daleifmodif}{
	\renewcommand\chapnamefont{\normalfont\Large\scshape\raggedleft\so} 
	\renewcommand\chaptitlefont{\normalfont\Large\bfseries\scshape} 
	\renewcommand\chapternamenum{} \renewcommand\printchaptername{} 
	\renewcommand\printchapternum{\null\hfill\feline@chm[2.5cm]\par} 
	\renewcommand\afterchapternum{\par\vskip\midchapskip} 
	\renewcommand\printchaptertitle[1]{\color{gray}\chaptitlefont\raggedleft ##1\par}
} 
\makeatother 
\chapterstyle{daleifmodif}
%
% UoB guidelines:
%
% The pages should be numbered consecutively at the bottom centre of the
% page.
\makepagestyle{myvf} 
\makeoddfoot{myvf}{}{\thepage}{} 
\makeevenfoot{myvf}{}{\thepage}{} 
\makeheadrule{myvf}{\textwidth}{\normalrulethickness} 
\makeevenhead{myvf}{\small\textsc{\leftmark}}{}{} 
\makeoddhead{myvf}{}{}{\small\textsc{\rightmark}}
\pagestyle{myvf}
%
% Oscar's command (it works):
% Fills blank pages until next odd-numbered page. Used to emulate single-sided
% frontmatter. This will work for title, abstract and declaration. Though the
% contents sections will each start on an odd-numbered page they will
% spill over onto the even-numbered pages if extending beyond one page
% (hopefully, this is ok).
\newcommand{\clearemptydoublepage}{\newpage{\thispagestyle{empty}\cleardoublepage}}
%
%
% Creates indexes for Table of Contents, List of Figures, List of Tables and Index
\makeindex
% \printglossaries below creates a list of abbreviations. \gls and related
% commands are then used throughout the text, so that latex can automatically
% keep track of which abbreviations have already been defined in the text.
%
% The import command enables each chapter tex file to use relative paths when
% accessing supplementary files. For example, to include
% chapters/brewing/images/figure1.png from chapters/brewing/brewing.tex we can
% use
% \includegraphics{images/figure1}
% instead of
% \includegraphics{chapters/brewing/images/figure1}
\usepackage{import}

% Add other packages needed for chapters here. For example:
\usepackage{lipsum}					%Needed to create dummy text
\usepackage{amsfonts} 					%Calls Amer. Math. Soc. (AMS) fonts
\usepackage[centertags]{amsmath}			%Writes maths centred down
\usepackage{stmaryrd}					%New AMS symbols
\usepackage{amssymb}					%Calls AMS symbols
\usepackage{amsthm}					%Calls AMS theorem environment
\usepackage{newlfont}					%Helpful package for fonts and symbols
\usepackage{layouts}					%Layout diagrams
\usepackage{graphicx}					%Calls figure environment
\usepackage{longtable,rotating}			%Long tab environments including rotation. 
\usepackage[applemac]{inputenc}			%Needed to encode non-english characters 
									%directly for mac
\usepackage{colortbl}					%Makes coloured tables
\usepackage{wasysym}					%More math symbols
\usepackage{mathrsfs}					%Even more math symbols
\usepackage{float}						%Helps to place figures, tables, etc. 
\usepackage{verbatim}					%Permits pre-formated text insertion
\usepackage{upgreek }					%Calls other kind of greek alphabet
\usepackage{latexsym}					%Extra symbols
\usepackage[super,square,numbers,
		     sort&compress]{natbib}		%Calls bibliography commands 
\usepackage{url}						%Supports url commands
\usepackage{etex}						%eTeXÕs extended support for counters
\usepackage{fixltx2e}					%Eliminates some in felicities of the 
									%original LaTeX kernel
\usepackage[spanish,english]{babel}		%For languages characters and hyphenation
\usepackage{color}                    				%Creates coloured text and background
\usepackage[colorlinks=true,
		     allcolors=black]{hyperref}              %Creates hyperlinks in cross references
\usepackage{memhfixc}					%Must be used on memoir document 
									%class after hyperref
\usepackage{enumerate}					%For enumeration counter
\usepackage{footnote}					%For footnotes
\usepackage{microtype}					%Makes pdf look better.
\usepackage{rotfloat}					%For rotating and float environments as tables, 
									%figures, etc. 
\usepackage{alltt}						%LaTeX commands are not disabled in 
									%verbatim-like environment
\usepackage[version=0.96]{pgf}			%PGF/TikZ is a tandem of languages for producing vector graphics from a 
\usepackage{tikz}						%geometric/algebraic description.
\usetikzlibrary{arrows,shapes,snakes,
		       automata,backgrounds,
		       petri,topaths}				%To use diverse features from tikz		
%		

% Todo package
\usepackage[colorinlistoftodos,prependcaption,textsize=tiny]{todonotes}
%

% Epigraph package
\usepackage{epigraph}
\setlength{\afterepigraphskip}{\baselineskip}
%\setlength\epigraphrule{0pt}
%

%Minted package/ Code formatting
\usepackage{xcolor}
\usepackage{minted}
%\usemintedstyle{colorful}	

\definecolor{mintedbackground}{rgb}{0.95,0.95,0.95}
\newcommand{\code}[1]{\colorbox{mintedbackground}{\texttt{#1}}}

\newcommand{\hminted}{\begin{minted}[
bgcolor=mintedbackground,
numberblanklines=true,
numbersep=10pt,
gobble=0,
frame=lines,
framerule=0.4pt,
framesep=4mm,
funcnamehighlighting=false,
tabsize=4,
obeytabs=false,
mathescape=false
samepage=false, %with this setting you can force the list to appear on the same page
showspaces=false,
showtabs =false,
texcl=false
]{haskell}}

\newcommand{\hmint}[1]{\mint[
bgcolor=mintedbackground,
numberblanklines=true,
numbersep=10pt,
gobble=0,
frame=lines,
framerule=0.4pt,
framesep=4mm,
funcnamehighlighting=false,
tabsize=4,
obeytabs=false,
mathescape=false
samepage=false, %with this setting you can force the list to appear on the same page
showspaces=false,
showtabs =false,
texcl=false
]{haskell}#1}

\newmintedfile[haskellcode]{haskell}{
bgcolor=mintedbackground,
linenos=true,
numberblanklines=true,
numbersep=10pt,
gobble=0,
frame=lines,
framerule=0.4pt,
framesep=4mm,
funcnamehighlighting=false,
tabsize=4,
obeytabs=false,
mathescape=false
samepage=false, %with this setting you can force the list to appear on the same page
showspaces=false,
showtabs =false,
texcl=false
}
				

%Reduce widows  (the last line of a paragraph at the start of a page) and orphans 
% (the first line of paragraph at the end of a page)
\widowpenalty=1000
\clubpenalty=1000
%
% New command definitions for my thesis
%
\newcommand{\keywords}[1]{\par\noindent{\small{\bf Keywords:} #1}} %Defines keywords small section
\newcommand{\parcial}[2]{\frac{\partial#1}{\partial#2}}                             %Defines a partial operator
\newcommand{\vectorr}[1]{\mathbf{#1}}                                                        %Defines a bold vector
\newcommand{\vecol}[2]{\left(                                                                         %Defines a column vector
	\begin{array}{c} 
		\displaystyle#1 \\
		\displaystyle#2
	\end{array}\right)}
\newcommand{\mados}[4]{\left(                                                                       %Defines a 2x2 matrix
	\begin{array}{cc}
		\displaystyle#1 &\displaystyle #2 \\
		\displaystyle#3 & \displaystyle#4
	\end{array}\right)}
\newcommand{\pgftextcircled}[1]{                                                                    %Defines encircled text
    \setbox0=\hbox{#1}%
    \dimen0\wd0%
    \divide\dimen0 by 2%
    \begin{tikzpicture}[baseline=(a.base)]%
        \useasboundingbox (-\the\dimen0,0pt) rectangle (\the\dimen0,1pt);
        \node[circle,draw,outer sep=0pt,inner sep=0.1ex] (a) {#1};
    \end{tikzpicture}
}
\newcommand{\range}[1]{\textnormal{range }#1}                                             %Defines range operator
\newcommand{\innerp}[2]{\left\langle#1,#2\right\rangle}                                 %Defines inner product
\newcommand{\prom}[1]{\left\langle#1\right\rangle}                                         %Defines average operator
\newcommand{\tra}[1]{\textnormal{tra} \: #1}                                                       %Defines trace operator
\newcommand{\sign}[1]{\textnormal{sign\,}#1}                                                   %Defines sign operator
\newcommand{\sech}[1]{\textnormal{sech} #1}                                                  %Defines sech
\newcommand{\diag}[1]{\textnormal{diag} #1}                                                    %Defines diag operator
\newcommand{\arcsech}[1]{\textnormal{arcsech} #1}                                       %Defines arcsech
\newcommand{\arctanh}[1]{\textnormal{arctanh} #1}                                         %Defines arctanh
%Change tombstone symbol
\newcommand{\blackged}{\hfill$\blacksquare$}
\newcommand{\whiteged}{\hfill$\square$}
\newcounter{proofcount}
\renewenvironment{proof}[1][\proofname.]{\par
 \ifnum \theproofcount>0 \pushQED{\whiteged} \else \pushQED{\blackged} \fi%
 \refstepcounter{proofcount}
 \normalfont 
 \trivlist
 \item[\hskip\labelsep
       \itshape
   {\bf\em #1}]\ignorespaces
}{%
 \addtocounter{proofcount}{-1}
 \popQED\endtrivlist
}
%
%
% New definition of square root:
% it renames \sqrt as \oldsqrt
\let\oldsqrt\sqrt
% it defines the new \sqrt in terms of the old one
\def\sqrt{\mathpalette\DHLhksqrt}
\def\DHLhksqrt#1#2{%
\setbox0=\hbox{$#1\oldsqrt{#2\,}$}\dimen0=\ht0
\advance\dimen0-0.2\ht0
\setbox2=\hbox{\vrule height\ht0 depth -\dimen0}%
{\box0\lower0.4pt\box2}}
%
% My caption style
\newcommand{\mycaption}[2][\@empty]{
	\captionnamefont{\scshape} 
	\changecaptionwidth
	\captionwidth{0.9\linewidth}
	\captiondelim{.\:} 
	\indentcaption{0.75cm}
	\captionstyle[\centering]{}
	\setlength{\belowcaptionskip}{10pt}
	\ifx \@empty#1 \caption{#2}\else \caption[#1]{#2}
}
%
% My subcaption style
\newcommand{\mysubcaption}[2][\@empty]{
	\subcaptionsize{\small}
	\hangsubcaption
	\subcaptionlabelfont{\rmfamily}
	\sidecapstyle{\raggedright}
	\setlength{\belowcaptionskip}{10pt}
	\ifx \@empty#1 \subcaption{#2}\else \subcaption[#1]{#2}
}
%
%An initial of the very first character of the content
\usepackage{lettrine}
\newcommand{\initial}[1]{%
	\lettrine[lines=3,lhang=0.33,nindent=0em]{
		\color{gray}
     		{\textsc{#1}}}{}}
%
% Theorem styles used in my thesis
%
\theoremstyle{plain}
\newtheorem{theo}{Theorem}[chapter]
\theoremstyle{plain}
\newtheorem{prop}{Proposition}[chapter]
\theoremstyle{plain}
\theoremstyle{definition}
\newtheorem{dfn}{Definition}[chapter]
\theoremstyle{plain}
\newtheorem{lema}{Lemma}[chapter]
\theoremstyle{plain}
\newtheorem{cor}{Corollary}[chapter]
\theoremstyle{plain}
\newtheorem{resu}{Result}[chapter]
%
% Hyphenation for some words
%
\hyphenation{res-pec-tively}
\hyphenation{mono-ti-ca-lly}
\hyphenation{hypo-the-sis}
\hyphenation{para-me-ters}
\hyphenation{sol-va-bi-li-ty}
%

%chapter subtitle
\newcommand\Chapter[2]{
  \chapter[#1]{#1:\\\Large\itshape#2}
}

%
\begin{document}
% UoB guidlines:
%
% Preliminary pages
% 
% The five preliminary pages must be the Title Page, Abstract, Dedication
% and Acknowledgements, Author's Declaration and Table of Contents.
% These should be single-sided.
% 
% Table of contents, list of tables and illustrative material
% 
% The table of contents must list, with page numbers, all chapters,
 % sections and subsections, the list of references, bibliography, list of
% abbreviations and appendices. The list of tables and illustrations
% should follow the table of contents, listing with page numbers the
% tables, photographs, diagrams, etc., in the order in which they appear
% in the text.
% 
\frontmatter
\pagenumbering{roman}
%
%
% File: Title.tex
% Author: V?ctor Bre?a-Medina
% Description: Contains the title page
%
% UoB guidelines:
% 
% At the top of the title page, within the margins, the dissertation should give the title and, if 
% necessary, sub-title and volume number. If the dissertation is in a language other than English, the 
% title must be given in that language and in English. The full name of the author should be in the 
% centre of the page. At the bottom centre should be the words ?A dissertation submitted to the 
% University of Bristol in accordance with the requirements for award of the degree of ? in the 
% Faculty of ...?, with the name of the school and month and year of submission. The word count of 
% the dissertation (text only) should be entered at the bottom right-hand side of the page.
%
%
\begin{titlingpage}
\begin{SingleSpace}
\calccentering{\unitlength} 
\begin{adjustwidth*}{\unitlength}{-\unitlength}
\vspace*{13mm}
\begin{center}
\rule[0.5ex]{\linewidth}{2pt}\vspace*{-\baselineskip}\vspace*{3.2pt}
\rule[0.5ex]{\linewidth}{1pt}\\[\baselineskip]
{\HUGE What Does It Mean To Be Reactive?}\\[4mm]
{\Large \textit{Crushing opinions with the power of Mathematics}}\\
\rule[0.5ex]{\linewidth}{1pt}\vspace*{-\baselineskip}\vspace{3.2pt}
\rule[0.5ex]{\linewidth}{2pt}\\
\vspace{6.5mm}
{\large By}\\
\vspace{6.5mm}
{\large\textsc{Eddy Bertoluzzo}}\\
\vspace{11mm}
\includegraphics[scale=0.6]{logos/TU_Delft_logo_RGB}\\
\vspace{6mm}
{\large Faculty of Electrical Engineering, Mathematics and Computer Science\\
\textsc{Delft University of Technology}}\\
\vspace{11mm}
\begin{minipage}{10cm}
A dissertation submitted to Delft University of Technology in accordance with the requirements of the \textsc{Master of Science} degree in the Faculty of Electrical Engineering, Mathematics and Computer Science.
\end{minipage}\\
\vspace{9mm}
{\large\textsc{October 2016}}
\vspace{12mm}
\end{center}
\begin{flushright}
%{\small Word count: ten thousand and four}
\end{flushright}
\end{adjustwidth*}
\end{SingleSpace}
\end{titlingpage}
\clearemptydoublepage
%
%
% File: abstract.tex
% Author: V?ctor Bre?a-Medina
% Description: Contains the text for thesis abstract
%
% UoB guidelines:
%
% Each copy must include an abstract or summary of the dissertation in not
% more than 300 words, on one side of A4, which should be single-spaced in a
% font size in the range 10 to 12. If the dissertation is in a language other
% than English, an abstract in that language and an abstract in English must
% be included.

\chapter*{Abstract}
\begin{SingleSpace}
\initial{W}hat does it mean to be Reactive? The concept of Reactive Programming has gained much traction in the last few years as a paradigm well-suited for the development of asynchronous  event-driven applications. Unfortunately, Reactive Programming has been at the center of much discussion, if not confusion, with regards to its definition, properties and identifying principles. In this work we are going to wield the most powerful tool available to software engineers, mathematics, in order to formally derive the reactive types and bring clarity to this much opinionated topic.
\end{SingleSpace}
\clearpage
\clearemptydoublepage
%
%
% file: dedication.tex
% author: V?ctor Bre?a-Medina
% description: Contains the text for thesis dedication
%

\chapter*{Dedication and acknowledgements}
\begin{SingleSpace}
\initial{H}ere goes the dedication.
\end{SingleSpace}
\clearpage
\clearemptydoublepage
%
%%
% File: declaration.tex
% Author: V?ctor Bre?a-Medina
% Description: Contains the declaration page
%
% UoB guidelines:
%
% Author's declaration
%
% I declare that the work in this dissertation was carried out in accordance
% with the requirements of the University's Regulations and Code of Practice
% for Research Degree Programmes and that it has not been submitted for any
% other academic award. Except where indicated by specific reference in the
% text, the work is the candidate's own work. Work done in collaboration with,
% or with the assistance of, others, is indicated as such. Any views expressed
% in the dissertation are those of the author.
%
% SIGNED: .............................................................
% DATE:..........................
%
\chapter*{Author's declaration}
\begin{SingleSpace}
\begin{quote}
\initial{I} declare that the work in this dissertation was carried out in accordance with the requirements of  the University's Regulations and Code of Practice for Research Degree Programmes and that it  has not been submitted for any other academic award. Except where indicated by specific  reference in the text, the work is the candidate's own work. Work done in collaboration with, or with the assistance of, others, is indicated as such. Any views expressed in the dissertation are those of the author.

\vspace{1.5cm}
\noindent
\hspace{-0.75cm}\textsc{SIGNED: .................................................... DATE: ..........................................}
\end{quote}
\end{SingleSpace}
\clearpage
\clearemptydoublepage
%
\renewcommand{\contentsname}{Table of Contents}
\maxtocdepth{subsection}
\tableofcontents*
\addtocontents{toc}{\par\nobreak \mbox{}\hfill{\bf Page}\par\nobreak}
\clearemptydoublepage
%
%\listoftables
%\addtocontents{lot}{\par\nobreak\textbf{{\scshape Table} \hfill Page}\par\nobreak}
%\clearemptydoublepage
%
%\listoffigures
%\addtocontents{lof}{\par\nobreak\textbf{{\scshape Figure} \hfill Page}\par\nobreak}
%\clearemptydoublepage
%
%
% The bulk of the document is delegated to these chapter files in
% subdirectories.
\mainmatter


\let\textcircled=\pgftextcircled
\chapter*[Introduction]{Introduction}
\label{chap:intro}
\addcontentsline{toc}{chapter}{Introduction}

\epigraph{\hspace{4ex}\textit{"Lasciate ogne speranza, voi ch'intrate."}}{--- Dante Alighieri,\\ \textit{Divina Commedia}}

 
With the evolution of technologies brought in by the new millennium and the exponential growth of Internet-based services targeting millions of users all over the world, the Software Engineers' community has been continuously tested by an ever growing number of challenges related to management of increasingly large amounts of user data\cite{furht2010handbook}. 

This phenomena is commonly referred to as Big Data. A very popular 2001 research report\cite{laney20013d} by analyst Doug Laney, proposes a definition of big data based on its three defining characteristics:

\begin{itemize}
\item \textit{Volume}: the quantity of data applications have to deal with, ranging from small - e.g. locally maintained Databases - to large - e.g. distributed File Systems replicated among data centers.
\item \textit{Variety}: the type and structure of data, ranging from classic SQL-structured data sets to more diversified and unstructured ones such as text, images, audio and  video. 
\item \textit{Velocity}: the speed at which data is generated, establishing the difference between pull-based systems, where data is synchronously pulled by the consumer, and push-based systems, more suited for handling real-time data by asynchronously pushing it to its clients.
\end{itemize}

Each of these traits directly influences the way programming languages, APIs and databases are designed today. The increasing volume calls for a declarative approach to data handling as opposed to an imperative one, resulting in the developer's focus shifting from how to compute something to what it is to be computed in the first place\cite{fahland2009declarative}. The diversification of data, on the other hand, is the main drive for the research and development of noSQL approaches to data storage. Lastly, the increase in velocity fuels the need for event-driven, push-based models of computation that can better manage the high throughput of incoming data\cite{meijer2012your}. 

In this context, the concept of \textit{reactive programming} has gained much traction in the developer's community as a paradigm well-suited for the development of asynchronous  event-driven applications\cite{bainomugisha2013survey}. Unfortunately, reactive programming has been at the center of much discussion, if not confusion, with regards to its definition, properties and principles that identify it\cite{meijer2014reactive}.

The goal of our work is to use mathematics as a tool to analyze today's commercial reactive libraries and understand how they relate to the theoretical concept of reactive programming.  We are going to do so by utilizing concepts and ideas from functional programming and category theory with the purpose of deriving the reactive types, and will subsequently continue with the development of a reference reactive library which builds upon the previously derived theoretical foundations. 

\section*{Contributions}
\addcontentsline{toc}{section}{Contributions}

To the best of our knowledge, we are not aware of any previous work which analyses reactive programming from a theoretical standpoint or derives its types and interfaces though the use of mathematics.

The most commonly known attempt at defining reactive programming is the Reactive Manifesto\cite{reactive-manifesto}, a document that tries to describe reactive systems in terms of design principles and conceptual architecture. Although certainly insightful, we feel the document targets a more managing-focused audience as opposed to software developers and engineers. 

Our work, on the other hand, aims at providing types and interfaces that describe the real essence of the reactive paradigm, aiding engineers that wish to use or develop reactive libraries in understanding and taking more informed decisions on the matter.

This goal is achieved by providing a mathematical derivation of the reactive types, starting from they interactive counterparts and making use of theoretical concepts from category theory. These derived types are then used in the implementation of a formal reactive library where the purpose is showing how the theoretical definitions given to the various components can effectively be translated into working code. 

Together with this formal definition of the paradigm, this work presents a reference implementation for a production level reactive library, as well as an highlight of the issues and challenges encountered when bridging from the theoretical foundations of reactive programming to a concrete implementation of a reactive API. 

With the help of this report and the associated code repository, any software engineer interested in the topic should be able to understand the theoretical foundations behind the reactive paradigm and develop a reactive library in any language of choice.

\section*{Overview}
\addcontentsline{toc}{section}{Overview}

Chapter 1 introduces the scope of our research, providing a definition of reactive programming, the motivation and reasoning behinf our research and an overview of the current technologies and APIs that claim to belong to the world of reactive programming. Chapter 2 presents the mathematical derivation of the reactive types and interfaces, starting from the definition of \code{Iterable} and ending with that of \code{Observable}. Chapter 3 builds the bridge between the formal definition of the reactive types and a production level implementation of the paradigm, highlighting the technical issues as well as analyzing the relations with the previously discussed formal definitions. Chapter 4 concludes with final thought and future work.

\section*{Notation \& Conventions}
\addcontentsline{toc}{section}{Notation \& Conventions}

In the exposition of our work we will make use of Haskell as the reference language. This decision is motivated by the language's strong connection with mathematics as well as it's clean syntax.

All the code presented in this report, a minimal complete theoretical implementation and a full blown library implementation of a reactive library can be found at the associated code repository on \href{https://github.com/Widar91/Thesis/tree/master/reactive-formalism/rx-hs}{Github}. 

\clearemptydoublepage

%\import{chapters/chapter01/}{chap01.tex}
\let\textcircled=\pgftextcircled
\chapter{Reactive Programming}
\label{chap:reactiveprogramming}

\epigraph{\hspace{4ex}\textit{The cold winds are rising in the North... Brace yourselves, winter is coming.}}{--- George R.R. Martin,\\ \textit{A Game of Thrones}}


In this chapter we are going to introduce the concept of reactive programming and motivate it's importance and relevance with regards to modern applications and the type of problems developers have to face nowadays. We are then going to introduce the most popular commercial libraries that claim to solve the reactive problem, with the purpose of giving the reader some context for our discussion and motivating the need for a mathematical formalization that abstracts over the class of problems these implementations set out to address.

\section{The Essence of Reactive Programs}

The use of the term reactive program in scientific literature is dated back to the mid-sixties\cite{scopus-reactive}. A relevant and insightful definition was given by G. Berry in 1991\cite{berry1991reactive} as he describes reactive programs in relation to their dual counterparts, interactive programs:

\begin{quote}
\hspace{4ex}``\textit{Interactive programs} interact at their own speed with users or with other programs; from a user point of view, a time-sharing system is interactive. \textit{Reactive programs} also maintain a continuous interaction with their environment, but at a speed which is determined by the environment, not by the program itself.''
\end{quote}

Interactive programs concretize the idea of a pull-based model of computation, where the program - the consumer in this case - has control over the speed at which data will be requested and handled. A perfect example of an interactive program is a control-flow structure such as a for-loop iterating over a collection of data: the program is in control of the speed at which data is retrieved from the containing collection and will request the next element only after it is done handling the current one.

Reactive programs, on the contrary, embody the idea of a push-based - or event-driven - model of computation, where the speed at which the program interacts with the environment is determined by the environment rather than the program itself. In other words, it is now the producer of the data - i.e. the environment - who determines the speed at which events will occur whilst the program's role reduces to that of a silent observer that will react upon receiving events. Standard example of such systems are GUI applications dealing with various events originating from user input - e.g. mouse clicks, keyboard button presses - and programs dealing with stock market, social media or any other kind of asynchronous updates.  

\section{Why Reactive Programming Matters}
\label{sec:whyrpmatters}

Considering the definition and examples of reactive programs we analyzed in the previous section, let's now try to formalize the class of problems the reactive programming paradigm is specifically well-suited for.

The table below provides a collection of types offered by common programming languages for handling data, parameterized over two variables: the size of the data, either one or multiple values, and the way data is handled, either by synchronous or asynchronous computations\cite{meijer2015spicing}.

\begin{center}
    \begin{tabular}{| l | l | l |}
    \hline
    & \textbf{One} & \textbf{Many} \\ 
    \hline
	\textbf{Sync} & \code{a} & \code{Iterable a} \\ 
	\hline
	\textbf{Async} & \code{Future a} & \textit{Reactive Programming} \\ 
	\hline
    \end{tabular}
    %\caption{The four fundamental effects}
\end{center}

The first row shows that synchronous functions come in two flavors: classic functions that return a single value of type \code{a} and functions that produce a collection of results of type \code{a}, abstracted through the \code{Iterable a} interface (See section \ref{sec:iterables}). These types of functions embody the standard imperative, pull-based approach to programming, where a call to a function/method synchronously blocks until a result is produced. 

Moving on to the second row, we encounter \code{Future a}, an interface representing an asynchronous computation that, at a certain point in the future, will result in a value of type \code{a}. Futures are generally created by supplying two callbacks together with the asynchronous computation, one to be executed in case of success and the other one in case of error. 

Programming languages, however, are not as well equipped when it comes to handling asynchronous computations resulting in multiple values - i.e. push-based collections. The issue lies in the fact that the program's control flow is dictated by the environment rather than the program itself - i.e. inversion of control -, making it very hard to model such problems with commonly known control structures, which are optimized for sequential models of computation. Traditional solutions typically involve developers manually trying to compose callbacks by explicitly writing CPS (continuation passing style) code\cite{meijer2015spicing}, resulting in what it's commonly referred to as \textit{Callback Hell}\cite{edwards2009coherent}.

The aforementioned class of problems reflects the definition of reactive programs we analyzed in the previous section, where the environment asynchronously - i.e. at its own speed - pushes multiple events to the program. The reactive programming paradigm sets out to provide interfaces and abstractions to facilitate the modeling of such problems as push-based collections. 

However, interfaces are only as good as the implementations that back them up. In the next section we are going to discuss and analyze various APIs and libraries that claim to embody the reactive paradigm, motivating our need for a mathematical formalization to aid in unifying these different approaches under a single set of interfaces, eliminating doubts and confusion when it comes to the essence of reactive programming. 

\section{Reactive Programming in the Real World}

Talk about FRP, Reactive Streams, Rx and other stuff. CIte survey paper
Claim that Observable is dual of Iterator in order to justify the why.

\subsection{Reactive Extensions}
\label{subsec:rx}

\subsection{Reactive Streams}
Try to define itself as reactive programming, mathematics prove the claim wrong, they implement asynchronous iterables. This does not make them useless, yet they are not a good solution for the problem we are trying to solve.

\todo[inline]{Probably put the derivation of Reactive Streams in the appendix, it's not really the focus, I guess.}

\subsection{Functional Reactive Programming}
\label{frp}
Much researched topic, beautiful in theory, yet based on continuous time, un-achievable with computers where time is inherently discrete. Connel Elliot, the inventor of FRP, says that commercial implementations are far from the original theory and do not really represent the essence of FRP
\clearemptydoublepage

\let\textcircled=\pgftextcircled
\Chapter{Into the Rabbit Hole}{Deriving the Observable}

\epigraph{\hspace{4ex}\textit{"It was much pleasanter at home," thought poor Alice, "when one wasn't always growing larger and smaller, and being ordered about by mice and rabbits. I almost wish I hadn't gone down the rabbit-hole -- and yet -- and yet -- ..."}}{--- Lewis Carol,\\ \textit{Alice in Wonderland}}


In this chapter we are going to derive the \code{Observable} interface starting from its dual counterpart, the \code{Iterable}, which, as we have seen in Chapter \ref{chap:reactiveprogramming}, represents embodies the idea of a pull-based model of computation and is the commonly adopted solution to dealing with synchronous computations resulting in multiple values.

\section{Iterables}
\label{sec:iterables}

An \code{Iterable} is a programming construct which enables the user to traverse a collection of data, abstracting over the underlying implementation\cite{gamma1995design}.

The interface and semantics of \code{Iterable}s were first introduced by the Gang of Four though their Iterator pattern\cite{gamma1995design}; today's most used programming languages expose the \code{Iterable} as the root interface in standard Collections APIs. 

The \code{Iterable} interface is generally fixed across programming languages, with the exception of naming conventions - C\# and related languages call it IEnumerable - and slight differences in the types, as we can see from the following definitions.\\
\todo{Should I use a datatype instead of class?}

\haskellcode{src/iterable_interface.hs}\\

Although the essence of the pattern is preserved in both definitions, we claim that the C\# version more clearly and accurately reflects the way side effects play a role in the usage of the interface: \code{moveNext} contains all the side effects of walking down the underlying collection and retrieving the next value while \code{current} can inspect the retrieved value multiple times in a pure way. The Java version, on the other hand, embeds the side effect in the \code{next} function, making it impossible to inspect the current value multiple times. For this reason and without loss of generality, we will make use of the C\# definition - modulo naming conventions - in the reminder of the discussion.\\
\todo{Mention Subscription here?}

\haskellcode{src/iterable_interface_final.hs}\\

The derivation that follows will require the use of a number mathematical concepts such as categorical duality, continuations, (co)products, (un)currying, covariance, contravariance and functors. We suggest the reader to get familiar with these topics before diving into the derivation. An accessible introduction to each can be found in Appendix \ref{app:a}.

\section{The Essence of Iterables}

The first step in deriving the \code{Observable} is to simplify our \code{Iterable} definition to a type that reflects its very essence; we are gonna do this by stripping the interface presented in the previous section of all the unnecessary operational features that only clutter our definition.

Let's start by taking a closer look at the \code{Iterator} interface; we can observe that the definition of the functions \code{moveNext} and \code{current} is equivalent to a single function which returns either a value - analogous to a \code{moveNext} call returning true and a subsequent invocation to \code{current} - or nothing - analogous to a call to \code{moveNext} returning false. 

Before we formalize this observation with a proper type, let us notice another effect that is hidden in the current definition of \code{moveNext}, but not made explicit by the its type: the possibility for an exception to be thrown by the function's body. 

By merging these considerations with the notion of coproducts and Haskell's \code{Either} and \code{Maybe} type, we obtain the following definition.\\

\haskellcode{src/iterable_coproducts_maybe.hs}\\

Note how, theoretically, \code{getIterator} could also throw an exception. We assume here, without loss of generality, that the function will never throw and will always be able to return an \code{Iterator} instance.

The next step is to forget about class instances and express our interfaces as simple types.\\

\haskellcode[firstline=1,lastline=2]{src/iterable_types.hs}\\

At this point, we want to put aside the operational concerns regarding exceptions and termination and assume the \code{Iterator} function will always return a value of type \code{a}.

\haskellcode[firstline=4,lastline=5]{src/iterable_types.hs}\\

We have now reached a point where no simplification is possible anymore. The obtained types reflect the essence of the Iterator patter: an \code{Iterable} is simply a function which, when invoked, produces an \code{Iterator} and an \code{Iterator} is itself a function producing a value of type \code{a} as a side effect. 

\todo{Reformulate this part better, ask Erik for advice.}
These types present some interesting properties; let's start by analyzing the \code{Iterator}, a function which, given nothing, will produce a value of type \code{a}. This type should sound familiar to the reader acquainted with object oriented programming as it precisely describes the notion of a getter function, i.e. a lazy producer of values. When looking at the relation between the \code{Iterator} type and its base component, \code{a}, we can observe how they are bound by a covariance relation:
\todo{vending machine example?}
\begin{displaymath}
\frac{A <: B}{() \rightarrow A <: () \rightarrow B} 
\end{displaymath}

The \code{Iterable}, on the other hand, is nothing more than a getter of another getter, the \code{Iterator} and therefore abides by the covariance relation as well.

To formally prove this intuition of a covariant relation, we instantiate the \code{Iterable}/\code{Iterator} types to a covariant \code{Functor}.\\

\haskellcode[firstline=3]{src/iterable_functor.hs}\\

We will see in the next section how these concepts are relevant in expressing and motivating the relevance of the duality between Iterables and Observables.

\section{Applying Duality}
\label{sec:applyingduality}

By now, the reader should be somehow familiar with the concept of duality, as it has appeared many times throughout our discussion, in concepts such as pull and push models of computation or interactive and reactive programs. Duality is, in fact, a very important general theme that has manifestations in almost every area of mathematics\cite{gowers2010princeton} (See Appendice \ref{app:a} for an introductory discussion on the topic). 

Starting from the fact that the \code{Iterable} interface embodies the idea of interactive programming, let's use the principle of duality to derive the \code{Observable} interface and see how it relates to the concept of reactive programming. In practice, this translates to the simple task of flipping the function arrows in the \code{Iterable} interface, taking us from a function resulting in a value of type \code{a} to one accepting an \code{a}.\\

\haskellcode{src/iter_to_obs.hs}\\

Not how the side effects are bound to function application rather than values, hence their flipped position in the \code{Observable} type.

The reader acquainted with functional programming will easily see the resemblance between the \code{Observable} type and a CPS function (See Appendice \ref{app:a}).\\

\haskellcode{src/obs_cont.hs}\\

It is clear how an \code{Observable} is nothing more than a CPS function where the result type \code{r} is instantiated to \code{IO ()}. To convince ourselves of this equivalence, let's think about the definition of a CPS function, i.e. a suspended computation which, given another function - the continuation - as argument, will produce its final result. This definition suits perfectly the idea behind \code{Observable} discussed in Section \ref{subsec:rx}: a function which will do nothing - i.e. is suspended - until it is subscribed to by an \code{Observer}.

A continuation, on the other hand, represents the future of the computation, a function from an intermediate result to the final result\cite{newbern-monads}; in the context of \code{Observable}s, the continuation represents the \code{Observer}, a function specifying what will happen to a value produced by the \code{Observable}, whenever it will become available, that is, whenever it will be pushed into the \code{Observer}. Since a continuation can be called multiple times within the surrounding CPS context, it is easy to see how this mathematical concept allows us to easily deal with multiple values produced at different times in the future.

In the previous section we have discussed many properties associated with \code{Iterable}s. Let's analyze now how these properties translate under dualisation and how they affect our new derived interface, the \code{Observable}.

First, moving from the observation that an \code{Iterable} is a getter of a getter, we can observe that the \code{Observable} plays exactly the opposite role, that is, a setter of a setter. The type \code{Observer :: a -> IO ()} represents, in fact, the essence of a setter function, whereas the \code{Observable} consists in nothing more than the simple task of applying the observer function to itself, producing a setter of setters. 

\todo{Probably will need to riformulate this.}
While the discussion about \code{Iterable}'s covariance was quite intuitive, things get a little bit more complicated when analyzing \code{Observable}s. Let's start by informally introducing the notion of positivity and negativity of types: we can interpret a function of type \code{f :: a -> b} as a way for us to produce a value of type \code{b}. In this context, \code{b} is considered to be positive with respect to the type \code{a -> b}. On the other hand, in order to apply the function, we are going to need a value of type \code{a}, which we will need to get from somewhere else; \code{a} is therefore considered to be negative w.r.t. the function type, as the function introduces a need for this value in order to produce a result. The point of this distinction is that positive type variables introduce a covariant relation between base and function type whereas negative type variables introduce a contravariant relation.

Analyzing \code{Iterable} within this framework is easy, the \code{Iterator} function contains a single type parameter found in a positive position, therefore resulting in a covariant relation; being the \code{Iterable} the result of applying the \code{Iterator} function to itself, we again result in a covariant relation w.r.t. the type parameter \code{a}.

The \code{Observer} function, on the contrary, introduces a need for a value of type \code{a}, resulting in a contravariant relation w.r.t. \code{a}. Again, the \code{Observable} function is the result of applying \code{Observer} to itself; surprisingly, this results in \code{a} being in a positive position. The intuition is easily understood by thinking about the rules of arithmetic multiplication: \code{a} is in negative position w.r.t. the \code{Observer} function, whereas the \code{Observer} is in negative position w.r.t. the \code{Observable}. This leads to \code{a} being negated twice, ultimately resulting in a positive position within the \code{Observable} function. Before we formalize this intuition, let's convince ourselves that \code{Observable}s effectively produces a value of type \code{a} by looking at an example:\\

\haskellcode{src/polarity.hs}\\

It is clear from this implementation that \code{randomValueObs} indeed produces a value of type \code{Int}, whereas the \code{Observer} introduces a need for such value in order to be applied. For more details on the positivity and negativity of functions and type variables, see \cite{pos-neg}\cite{pos-neg2}.

Just as we did with \code{Iterable} we can formally prove the covariant and contravariant relations by instantiating \code{Observable} to a \code{Functor} instance and \code{Observer} to a Cofunctor one (Haskell uses the name \code{Contravariant} instead).\\

\haskellcode{src/observable_functor.hs}

\section{Observables are Continuations}
\label{obscont}

In the last section we mentioned how the \code{Observable} interface is equivalent to a CPS function. We are now going to formally prove this claim by providing a basic implementation of \code{Observable} as a type alias of Haskell's continuation monad. The Haskell language provides a monad construct for expressing continuations, as well as a monad transformer which allows the user to stack the functionality of continuations on top of other monads. A monad transformer is exactly what we need in order to express our \code{Observable} function producing a result in the IO monad.\\

\haskellcode[firstline=3,lastline=14]{src/observable_contT.hs}\\

At this point we have all the necessary tools to create and run an \code{Observable}.\\

\haskellcode[firstline=16]{src/observable_contT.hs}\\

Notice how an observable only pushes values once subscribed to and acts as a suspended computation otherwise. 

The code abobe, being a toy example fails to show some fundamental properties associated with this new interface; in particular, it fails to show how \code{Observable}s can actually handle asynchronous sources of data. The following snippet of code contains a more realistic and meaningful example of an \code{Observable} handling keyboard presses, asynchronous events by nature. It is worth noticing how our basic implementation of \code{Observable} based on continuations works just as well as a full blown one in terms of its core capability of handling asynchronous data.\\ 

\todo[inline]{Put the actual demo that listens to keyboard presses. In the example, press 3 times and see that all the events are handled and none are lost.}
\haskellcode[firstline=13]{src/create_run_demo.hs}\\

This example brings us to the following observation: \code{Observable}s are capable of handling asynchronous data sources, yet the means by which the data is handled are not asynchronous by default. This is a common misconception and source of much confusion among the community: an \code{Observable}, by default, synchronously handles it's incoming data, blocking the next incoming events whilst processing the current one. This behavior is not fixed though: as we will see in Chapter \ref{chap:outoftherabbithole}, it is possible to make use of \code{Scheduler}s in order to alter the control flow of the data processing and dispatch the work to other threads.

At this point in the discussion we have arrived to a working implementation of a push based collection purely derived from mathematical and categorical concepts such as duality and continuations. In the next Chapter we will take the necessary steps to transform this basic implementation into a full blown API that can be used in real world applications.
\clearemptydoublepage

\let\textcircled=\pgftextcircled
\Chapter{Out of the rabbit hole}{Towards a usable API}
\label{chap:outoftherabbithole}

\epigraph{\hspace{4ex}\textit{
''In theory there is no difference between theory and practice; in practice there is.''}}{--- Nassim Nicholas Taleb,\\\textit{Antifragile - Things that Gain From Disorder}}

So far we focused our analysis on the essence of the \code{Observable} interface, setting aside the many operational concerns that would come up when trying to implement these concepts into a usable, commercial API. In this chapter we are going to build the bridge between our theoretical definition of \code{Observable} and a concrete and usable implementation of a reactive library, to which we will refer to as Rx.

The goal of this chapter is to provide the reader with a reference implementation of a reactive library as well as to highlight the challenges and issues that emerge when trying to build the bridge between theory and practice. The implementation choices presented below are in no way prescriptive, instead, they aim to describe the problem in the most clear way, in order to stimulate awareness rather than blindly guide the reader to a solution. \todo{Say something about why I chose to go with the RxJava style of implementation?}

In the remainder of the discussion, we are going to introduce the \textit{Reactive Contract}, a set of assumptions on the reactive types our library is going to build upon, \textit{Schedulers}, which will allow us to bring concurrency into our reactive equation, \textit{Subscriptions}, used to implement a mechanism for premature stream cancellation and finally, \textit{Operators}, the means with which we will make our reactive streams composable.

For the sake of clarity and completeness, the following set of interfaces represents the starting point for our discussion:\footnote{\code{subscribe} has been renamed to \code{\_subscribe} in order to avoid naming conflicts later on in the discussion and reflect the fact that it should not be used directly by the user.}

\hminted 
newtype Observable a = Observable 
	{ _subscribe :: Observer a -> IO ()
	}
data Observer a = Observer 
	{ onNext      :: a -> IO ()
	, onError     :: SomeException -> IO ()
	, onCompleted :: IO ()
	}
\end{minted}


It is worth noting that even though the \code{Observable}'s theoretical foundations lie in the realm of functional programming, the road to making it usable is full of obstacles that are often better tackled using imperative programming features, such as state. As much as I personally prefer a functional and pure approach to programming, I will favor, in the rest of the discussion, the solution that most clearly and easily solves the problem, be that functional or imperative. As mentioned previously, Appendix \ref{app:b} contains a reference implementation of a Reactive Library which implements the features presented in this chapter utilizing existing constructs from functional programming.

\section{The Reactive Contract}
\label{sec:contract}

The \code{Observable} and \code{Observer} interfaces are somewhat limited, in their expressive power, to only argument and return types of they functions. The reactive library we are going to build is going to make more assumptions than the ones expressible by the type system. Although limiting, in a sense, the freedom with which the reactive interfaces can be utilized, this set of assumptions - the \textit{Contract} - greatly facilitates reasoning about and proving correctness of reactive programs\cite{MS2010-RxDesign}.

In later sections, we will refer back to these assumptions when discussing the actual implementation of our reactive library.

\subsection{Reactive Grammar}
\label{ass-grammar}
The first assumption we are going to introduce involves restrictions on the emission protocol of an \code{Observable}. Events propagated to the \code{Observer} continuation will obey the following regular expression:

\begin{center}\code{onNext* (onError | onCompleted)?}\end{center}

This grammar allows streams to propagate any number - 0 or more - of events through the \code{onNext} function, optionally followed by a message indicating termination, be that natural - through \code{onCompleted} - or due to a failure - through \code{onError}. Note how the optional nature of a termination message allows for the existence of infinite streams. 

This assumption is of paramount importance as it guarantees that no events can follow a termination message, allowing the consumer to effectively determine when it is safe to perform resource cleanup operations. 

\subsection{Serialized Messages}
\label{ass-serialized}
Later in this chapter we will see how we can introduce concurrency in our reactive library through the use of the \code{Scheduler} interface. From a practical point of view, this means that it will be possible for different messages, to arrive to an \code{Observer} from different execution contexts. If all \code{Observer} instances would have to deal with this scenario, the code in our library would soon become cluttered with concurrency-related housekeeping, making it harder to maintain and reason about.

For this reason, we assume that messages will always arrive in a serialized fashion. As a consequence, operators that deal with events from different execution contexts - e.g. combiner operators - are required to internally perform serialization.

\subsection{Best Effort Cancellation}
\label{ass-besteffort}
The next assumption involves premature stream cancellation via \code{Subscription}s and the function \code{unsubscribe}, used in order to stop the observation of events from an \code{Observable}; we are going to assume that whenever \code{unsubscribe} is invoked, the library will make a best effort attempt to stop all the ongoing work happening in the background. The reason is simple: it is not always safe to abort work that is being processed - e.g. database writes. Although the library might still complete the execution of pending work, its results are guaranteed not to be propagated to any \code{Observer} that was previously unsubscribed.

\subsection{Resource Cleanup After Termination}
As we mentioned in assumption \ref{ass-grammar}, the guarantee that no events will occur after the first termination message makes it possible to determine when resource cleanup operations are safe to perform. We will now make one step further and assume that resources \textit{will} be cleaned up immediately after termination. This will make sure that any related side-effect will occur in a predictable manner.

\section{Concurrency with Schedulers}
\label{schedulers}

At the end of Section \ref{obscont} we discussed how \code{Observable}s, by default, handle data by means of a synchronous pipeline, blocking the processing of successive elements via the call stack. It is worth mentioning again how this synchronous processing does not affect the ability of \code{Observable}s to handle asynchronous data.

However, this synchronous behavior might not always be the best solution, especially in real world applications, where we might want to have a thread dedicated to listening to incoming events and one which processes them. Enter the \code{Scheduler} interface, an orthogonal\cite{wiki:orthogonality} structure w.r.t. \code{Observable} which allows us to introduce concurrency into our reactive equation. 

\code{Scheduler}s allow us to to alter the control flow of the data processing within an observable expression, introducing a way to dispatch the work of any number of operators to be executed within the specified context, e.g. a new thread.

The \code{Scheduler} interface looks like the following \footnote{The interface presented in this section is the result of a simplification of the actual one, which involves \code{Subscription}s. We will discuss the impact of \code{Subscription}s on \code{Schedulers} in the next section; suffices to know that the version presented here has no negative effects w.r.t the generality of our discussion.}.

\haskellcode[firstline=31,lastline=34]{src/scheduler.hs}

\code{Scheduler}s expose two functions which are essentially equal, modulo arbitrary delays in time. Both of these functions take an \code{IO} action as input and dispatch it to the appropriate execution context, producing a side effect. 

To better understand \code{Scheduler}s, let us present the implementation of one of them, the newThread scheduler, which allows us to dispatch actions to a new, dedicated thread.

\haskellcode[firstline=36,lastline=48]{src/scheduler.hs}

The \code{newThread} function gives us a side effectful way of creating a \code{Scheduler} by generating a new execution context - i.e. a new thread - and setting up the necessary tools for safe communication with it. The \code{Scheduler} functions we are provided, on the other hand, simply write the input \code{IO} action to the channel and return, effectively dispatching the execution of those actions to the new thread. 

Up to this point we haven't mentioned \code{Observable}s at all. This is the reason why we previously claimed that \code{Scheduler} and \code{Observable} are connected by an orthogonal relationship: the two interfaces are independent from one another, yet, when used together within an observable expression, they provide the user with greater expressive power w.r.t. concurrency. 

The only thing missing now is a way for us to combine the functionality of these two interfaces: \code{observeOn} and \code{subscribeOn} are the operators that will aid us on this task. The former will allow us to dispatch any call to an observer continuation on to the specified execution context, whereas the latter will allow us to control the concurrency of the \code{Observable} subscribe function.

For the sake of completeness and understandability, the following snippet contains a simple implementation of the \code{observeOn} operator together with a sample usage.

\haskellcode[firstline=50]{src/scheduler.hs}


\subsection{A Note on the Concept of Time}

Our discussion on push-based collections so far has not once mentioned the concept of time. This might appear strange, especially to the reader familiar with Functional Reactive Programming, where functions over continuous time are at the foundations of the theory. This dependency on continuous time comes at a great cost: commercial FRP libraries fail to successfully implement the concepts found in the theory\cite{elliott2014denotational} as they cannot avoid simulating continuous time and approximating functions operating over it, being this concept inherently discrete in the context of computers.

Rx, on the other hand, completely sheds the notion of time from the notion of reactivity\cite{meijer2010observable}, shifting its focus, with the help of \code{Scheduler}s, to concurrency instead. Time still plays a role, although indirect, within the library: events are processed in the order they happen, and operators make sure such order is maintained, ultimately handing over to the user a stream of time-ordered events.

\subsection{A Note on Orthogonality}

Previously we discussed how concurrency is an orthogonal concept w.r.t. Rx - i.e. introducing concurrency does not affect or pollute the definition of our reactive interfaces. This statement is only true from a abstract point of view, falling short of its promises when looking from an implementation perspective, in particular, when dealing with combiner operators (see Section \ref{operators}) such as \code{(>>=)} or \code{combineLatest}. These operators will not work at their full potential in a synchronous setting, due to the fact that subscribing to a stream will consume it entirely - or forever process, in the case of an infinite stream - before allowing the operator to subscribe to a different one, effectively making interleaving of events impossible. 

The problem is gracefully solved with the introduction of \code{Scheduler}s, which, by allowing for \code{Observable}s to be executed on different contexts, indirectly make it possible for interleaving to happen and for combiner operators to work at their full potential. This comes at a cost: combiners operators are required to perform message serialization (see assumption \ref{ass-serialized}) as well as internal state synchronization as, with the introduction of concurrency, messages and state changes can now originate from different execution contexts.

\section{Subscriptions}
\label{sec:subscriptions}

With schedulers, we are now able to handle observable streams from different execution contexts. The next step in making Rx ready for a production environment is to add a mechanism that will allow us to stop a stream from anywhere in our program, whenever we don't require it's data anymore - i.e. a mechanism that will allow the user to communicate to the \code{Observable} that one of it's \code{Observer}s is no longer interested in receiving its events. 

\todo[inline]{Is the following clear?}
We discussed in the previous section how schedulers effectively boost the expressive power of our reactive expressions by introducing concurrency and interleaving among events originating from different streams. Introducing a cancellation mechanism, on the other hand, is a purely practical concern: although very useful from a practical perspective, especially in the context of resource management, it doesn't impact expressive power from a \textit{reactive} point of view.

The means by which we are going to introduce a cancellation mechanism inside our reactive equation is though the \code{Subscription} datatype. From a functionality point of view, what we are aiming for is for the \code{\_subscribe} function to hand back a \code{Subscription} whenever invoked; users will later be able to use this \code{Subscription} in order to prematurely cease the observation of a stream. This design is closely related to Dispose pattern utilized in the .NET framework\cite{cwalina2008framework}.

The first step in designing a new feature is to understand how the already existing interfaces will be affected by the newly introduced one; starting from our informal definition of \code{Observable} from section \ref{obscont}, let's now define a new function \code{(\%)}, which incorporates the notion of returning a \code{Subscription} and see how this is going to affect our types:

\hminted
%a  :: a -> IO Subscription
%?a :: (Either SomeException (Maybe a) -> IO ()) -> IO Subscription
\end{minted}


With this change, each execution of the \code{Observable} function now returns a \code{Subscription}, a means for the user to prematurely terminate the processing of the stream. 

The next question is the following: to whom does a \code{Subscription} belong to? The key observation in addressing this question is that an \code{Observable} can be subscribed to by multiple \code{Observer}s; our goal is to provide a mechanism that will allow for a fine-grained control over which \code{Observer} is supposed to stop receiving events. The answer is then straightforward: the notion of subscription is tight to that of observer. The following snippet reflects this observation:

\hminted
$a  :: (Subscription, Either SomeException (Maybe a) -> IO ())
%$a :: (Subscription, Either SomeException (Maybe a) -> IO ()) 
	-> IO Subscription
\end{minted}


Let's quickly summarize what we have discussed so far: a subscription is some object which will allow us to prematurely stop observing a specific stream; since any stream can be subscribed to by multiple observers, we need to associate subscriptions to observers as opposed to observables. Lastly, a subscription is returned every time an observer is subscribed to a stream through the \code{\_subscribe} function. The following modifications to our reactive interfaces reflect these ideas:

\haskellcode[firstline=13,lastline=21]{src/subsciption.hs}

So far we have talked a lot about \code{Subscription}s, yet we haven't clarified what the type really looks like. The general idea is to have \code{Subscription} record the state of the \code{Observer} w.r.t. the \code{Observable} it is subscribed to - be that subscribed or unsubscribed. This can be easily achieved with a variable \code{\_isUnsubscribed :: IORef Bool} initialized to \code{False}, indicating that the associated \code{Observer} is initially not unsubscribed.

From a practical point of view, it is useful to augment \code{Subscription} with some additional functionality. The following code shows a definition of \code{Subscription} which incorporates an \code{IO ()} action to be executed at unsubscription time. This is particularly useful when we want to associate resource cleanup actions to the termination - be that forced or natural - of a stream observation. Additionally, it is useful to make the type recursive, allowing \code{Subscription}s to contain other values of the same type. This will be extremely useful for internal coordination of operators such as \code{(>>=) :: Monad m => m a -> (a -> m b) -> m b}, where each input value will spawn and subscribe a new \code{Observable}, whose subscription should be linked to the original one. Section \ref{operators} will extensively discuss this matter.

\hminted
data Subscription = Subscription
	{ _isUnsubscribed :: IORef Bool
	, onUnsubscribe   :: IO ()
	, subscriptions   :: IORef [Subscription]  
	}
\end{minted}


%We will see, later on, how we can further augment this interface in order to incorporate additional functionality useful for internal operator coordination as well as management of schedulers. 

It's now time to introduce the two functions at the hearth of the whole cancellation mechanism: \code{unsubscribe} will take care of modifying the state carried by the \code{Subscription} - i.e. setting \code{\_isUnsubscribed} to \code{True} -  as well as execute the associated \code{IO ()} action, whereas \code{subscribe} will simply act as a proxy for the original \code{\_subscribe} function from the \code{Observable} interface.

%It's now time to introduce the two functions at the hearth of the whole cancellation mechanism: \code{unsubscribe} will take care of modifying the state carried by the \code{Subscription} - i.e. setting \code{\_isUnsubscribed} to \code{True} - as well as execute the associated \code{IO ()} action, whereas \code{subscribe} will act as a proxy for the original \code{\_subscribe} function from \code{Observable}, calling it with a special \code{Observer} as an argument, which will prevent events from propagating to the downstream once the \code{Subscription} has been unsubscribed. Note how, for this mechanism to work, it is necessary for the two \code{Observer}s to share the same \code{Subscription}.

\haskellcode[firstline=71,lastline=84]{src/subsciption.hs}

The \code{safeObserver} utilized by the \code{subscribe} function is of crucial importance to the functionality of our library and its the reason why we need to proxy the original \code{\_subscribe} function: its implementation, in fact, embeds two of the reactive contract assumptions introduced previously. The safe onNext/onError/onCompleted functions implement the subscription mechanism, preventing, through the \code{ifSubscribed} function, events from propagating to the underlying \code{Observer}, once the related \code{Subscription} has been unsubscribed. By doing so, it is easy to see how assumption \ref{ass-besteffort} is satisfied: unsubscribing from a stream does not force the stop of any outstanding work, yet it is made sure that any result produced after unsubscribing, if any, will not be delivered to the downstream \code{Observer} - i.e. the \code{Observer} supplied by the user. Additionally, \code{safeObserver} allows the enforcement of the reactive grammar seen in assumption \ref{ass-grammar}; this is done by calling \code{unsubscribe} as soon as the first termination message - be that \code{onError} or \code{onCompleted} - arrives, effectively preventing any additional event from being propagated.

Note that, with the current implementation of the subscription mechanism, an \code{Observer} can only be subscribed once and only to a single \code{Observable}, as, once its \code{Subscription} is unsubscribed, the \code{\_isUnsubscribed} field is never reset to \code{False}. This convention is shared by many already existing implementations of reactive libraries such as the onces under the ReactiveX umbrella\cite{rxjava-wiki}.

Now that we have a clear idea of how the subscription mechanism is supposed to work and how it is integrated into our library, let's take a look at a few observations and concerns that involve it. 

\subsection{Impact on Schedulers}

In the previous sections we discussed a simplified version of the \code{Scheduler} interface that was glossing, without loss of generality, over details regarding \code{Subscription}s. In practice, it is useful to associate \code{Subscription}s not only to \code{Observer}s but to \code{Scheduler}s as well. 

\haskellcode[firstline=95,lastline=99]{src/scheduler_subscription.hs}

With this version of the interface, each scheduled action returns a \code{Subscription}, offering fine grained control over the actions to be executed; at the same time, a \code{Subscription} is also associated to the \code{Scheduler} as a whole, allowing the user to perform cleanup actions on the \code{Scheduler} itself once \code{unsubscribe} is called. This is best shown with a new example implementation of the newThread scheduler and observeOn operator:

\haskellcode[firstline=101,lastline=130]{src/scheduler_subscription.hs}

The code is mostly equal to the one presented in section \ref{schedulers}. The most relevant change can be found at line 107, where we create a subscription for the newThread scheduler with an action that simply kills the thread\footnote{Absolutely not safe, but it's good enough for the sake of our example.}. On line 122 we then add this subscription to the one carried by the downstream observer. In this way, unsubscribing from the downstream subscription will trigger a waterfall effect that will eventually unsubscribing the scheduler's one as well, effectively killing the thread associated to it.

On a last note regarding the relationship between schedulers and subscriptions, it is worth mentioning how the subscription mechanism only works in the presence of schedulers. As we mentioned before, in fact, Rx is synchronous by default in the processing of its data. This means that the program would return from the invocation of the \code{subscribe} function only after it has fully processed the stream, effectively rendering the subscription mechanism ineffective, as it would not be possible to invoke \code{unsubscribe} whilst the \code{Observable} is active. With the introduction of schedules and different execution contexts, this problem disappears and the mechanism works as intended.

\subsection{Formalizing Subscriptions}

In Section \ref{obscont} we saw how an \code{Observable}, at its essence, is nothing more than an instance of the Continuation Monad, where the three types of events that can occur are materialized into a single datatype, \code{Event a}, as opposed to being handled by three different continuations.

In the discussion that follows we are going to try and understand what the essence of the subscription mechanism is and how it relates to our formal definition of observable as a continuation. 

As we mentioned before, a subscription is strongly tight to the notion of observer, as an observable can be subscribed-to multiple times. Although, we can be more specific than this and notice that a subscription is actually tight to the execution of an observable. These two takes on subscriptions are effectively the same thing: an observer can only be subscribed a single time to a single observable, creating the unique link between the subscription and a single execution of the observable function. This perspective is very insightful, as it hints to the fact that a subscription should be immutable within the context of an observable execution. Another observation is that the subscription needs to be retrievable from an observable for a number of reasons, the most important of which being to check whether the subscriber is unsubscribed before pushing any additional events.

The properties of subscription that we just discussed are very similar to the idea of environment variables, shared by computations yet immutable in their nature. The Haskell programming language exposes a monad construct for such computations, the Reader - Environment - Monad: in the remained of this section, we are going to model the subscription mechanism as a Reader monad transformer on top of our previous definition of observable as a continuation.

\hminted
type Observer   a = Event a -> IO ()
type Observable a = ReaderT Subscription (ContT () IO) (Event a)

subscribe :: Observable a -> Observer a -> IO Subscription
subscribe obs obr = do 
    subscription <- emptySubscription
	safeObserver  = enforceContract obr
    runContT (runReaderT obs s) safeObserver
    return s
    
enforceContract :: Observer a -> Observer a
enforceContract obr = ...
\end{minted}


This formalization is very insightful under many points of view: first of all, in the same way as schedulers, it is completely orthogonal to the definition of observable we previously had: the original definition did not change, yet the mechanism was, in a way, glued on top of it. This observation becomes very clear when looking at line 8 in the above snippet: the \code{subscribe} function first runs the reader transformer, resulting in continuation monad that will have the environment variable available within its context. Notice how each call to subscribe will effectively create a new subscription and pair it to the execution of the observable. 

A natural question now is: why didn't we use this technique for implementing subscriptions in the "real world" implementation from the previous paragraph? There, we had to change the definitions of our interfaces, losing orthogonality as a consequence. The reason is simply clarity, the interfaces look more clear than reading readerT, especially if we want to use that implementation as a reference for other languages. On the other hand, the goal of this paragraph is focusing on the essence of the subscription mechanism, hence the use of theory-related constructs such as monads.

Notice how this formalization only focuses on augmenting the definition of observables with subscriptions. The actual logic of the mechanism remains unchanged and is abstracted away through the \code{enforceContract} function.

As a final point, it is worth noting that what we presented so far is obviously not the only way we can formalize the subscription mechanism. Many other definitions have been tried out during the course of this work, yet all of the others ended up granting too much or too less power to the resulting mechanism and were therefore discarded. Examples include:

\hminted
type Observable a = StateT Subscription (ContT () IO) (Event a)
type Observable a = Cont () (StateT Subscription IO) (Event a)
\end{minted}


%it would be perfect in order to thread Subscriptions throughout execution making it usable inside operators, since the call to the continuation would return a state which would then be sequenced by >>=. This method fails with schedulers, in particular newThread since the state of the action executed on the new thread would be disconnected from the threading mentioned above. The other thread would get a state but it would not know what to do with it and would not have any ways to connect it to the original one passed by the subscribe. Another way is to use mapStateT to map the IO action that will result from the state to an action on the other thread. Again, this method won't work, 

%On the motivation for a Subscription: it is needed so the user can cancel work at any time. This implies that a scheduler is used. If this is not the case the subscription will be returned synchronously after the execution of the whole stream. It will therefore be already unsubscribed and calling unsubscribe will be a NoOp. In the case that we actually use schedulers then we can unsubscribe from anywhere in our program. 

%Now the question is the following: can we reproduce the behaviour of unsubscribe with operators so that we can eliminate subscriptions altogether? That is, we can hide them to the outside and use them only inside the stream, we still need them but we don't necessarily need to return them when we subscribe. The behavious can be easily replaced by the use of takeUntil(Observable a) where we pass in a subject that will be signaled when we want to stop the stream. The takeUntil operator fires events from the upstream up until the point in which we signal the subject, then stops the stream and unsubscribes. 

%With this approach we seemingly lose one thing, i.e. the ability to specify what happens at unsubscription time; this functionality can simply be regained by a subscribe function that takes the unsubscribe action as a parameter and runs it when the stream is unsubscribed.

\section{Operators}
\label{operators}


%- how does it work with subscriptions?

%- formalizing operators: stateful objects, nothing to do about that. Drawing of continuations called

With schedulers and subscriptions we can handle asynchronous streams of data from different execution context and cease our observation at any point in time. The last feature before we can consider our reactive library ready for use by developers is the introduction of a set of higher order functions that, given one or more \code{Observable}s, will allow access to its underlying elements - providing ways to transform, compose and filter them - while abstracting over the enclosing data structure. These higher order functions, also referred to as operators, are a common technique when it comes to abstractions over data structures; examples can be found in many commonly known programming languages, where data structures such as iterables, lists, trees, sets, ..., offer a wide range of functions - map, filter, flatmap, concat, etc - that allow the user to operate on its internal elements without requiring any knowledge of the structure that contains them.

\subsection{Functor, Applicative, Monad}

Back in section \ref{sec:applyingduality}, we defined the Functor instance for the Observable type, effectively augmenting our reactive type with our first operator, \code{map :: (a -> b) -> Observable a -> Observable b}, allowing the user to transform streams of elements of type \code{a} into streams of elements of type \code{b} by applying the input function of type \code{a -> b} to each incoming element in the input stream.

The next operator we are going to introduce is \code{lift}. This operator takes a function defined on the Observer level and lifts it to a more general context, that of Observable\cite{lift}. This pattern is very common in the field of Functional Programming, and it is often used in order to provide an abstraction that facilitates accessing values nested inside container structures such as Functors, Applicatives or Monads. In the context of Observables, this function will result very useful as we will be able to define operators on the Observer level and later make them accessible in the context of Observables simply by lifting the operation.

\haskellcode[firstline=23,lastline=24]{src/subsciption.hs}

As an example, we re-propose here the implementation of the Functor instance for Observable, this time using the \code{lift} function:

\haskellcode[firstline=26,lastline=31]{src/subsciption.hs}

As we discussed in section \ref{sec:applyingduality}, the \code{Observable} type, at its essence, is an instance of Applicative Functor and Monad from a cathegory theory perspective. This observation motivates our next effort of trying and defining such instances on the latest version of our interface. We will not bother anymore with proving the associated laws, as it would be a very difficult task, given the complicated nature of the implementation of such functions which now needs to handle subscriptions as well as synchronized state due to schedulers.\todo{Say more about the fact that they dont respect the laws, but just like the try, that is not a monad, they are ok? Meaning try doesnt respect the laws but nobody really cares, it does what it has to do and from an interface perspective it acts as if it was a monad.}

\haskellcode[firstline=139,lastline=189]{src/subsciption.hs}

The function \code{pure} does nothing more than wrapping a value into an \code{Observable} whereas \code{(<*>)} applies the most recent function emitted by the first \code{Observable} to the most recent element emitted by the second one. The implementation utilizes the more general \code{combineLatest} operator, which allows to combine two streams into one by emitting an item whenever either of the two emits one - provided that each of them has emitted at least one. 

Last but not least comes the Monad instance for our Observable type:

\haskellcode[firstline=191,lastline=243]{src/subsciption.hs}

Note how this implementation if \code{(>>=)} shows the need for the introduction of children subscriptions that we discussed in section \ref{sec:subscriptions}: whenever the outer observable is unsubscribed, we want to automatically unsubscribe any observable that has previously been created by the function passed to \code{(>>=)}.

\subsection{The Essential Operators}
Operators can theoretically be infinite in number, as infinite are the transformations that can be done to elements of an observable stream. Practice shows, though, that a relatively small subset of operators and the composition of these, suffices to express the majority of the use cases encountered by users. Leveraging the power of composability of operators is advantageous towards the design of a simple yet powerful API.

Operators can be grouped into categories by looking at their characteristics; it is not the purpose of this work to list every possible operator and its semantics, yet, for the sake of completeness, table \ref{ess-oper} presents the most important categories\footnote{An implementation of these operators can be found at the repository associated with this work.} and the associated operators:\todo{Put operator description in the table?}

\begin{table}[]
\thisfloatpagestyle{empty}
\centering
\caption{The essential operators}
\label{ess-oper}
\begin{tabular}{ll}
                                                        &                                                                                                                                                                                  \\ \hline
%\multicolumn{1}{|l|}{\textbf{Creation}}                 & \multicolumn{1}{l|}{}                                                                                                                                                            \\ \hline
\multicolumn{1}{|l|}{\textbf{Transformation}}           & \multicolumn{1}{l|}{\begin{tabular}[c] {l}\\ buffer\\ bufferWithSkip\\ groupBy\\ fmap  \\ scan \\ scanLeft\\ sample\\ throttle\\ window\\{\color{white} a}\end{tabular}}                          \\ \hline
\multicolumn{1}{|l|}{\textbf{Filtering \& Conditional}} & \multicolumn{1}{l|}{\begin{tabular}[c]{l} \\ filter\\ distinctUntilChanged\\ skip\\ skipUntil\\ skipWhile\\ take\\ takeUntil\\ untilTake \\{\color{white} a} \end{tabular}}                         \\ \hline
\multicolumn{1}{|l|}{\textbf{Combining}}                & \multicolumn{1}{l|}{\begin{tabular}[c]{l} \\ (\textgreater\textgreater=)\\ (\textless*\textgreater)\\ concat\\ startsWith\\ withLatestFrom\\ zip\\ zipWithBuffer \\{\color{white} a} \end{tabular}} \\ \hline
\multicolumn{1}{|l|}{\textbf{Error Handling}}           & \multicolumn{1}{l|}{\begin{tabular}[c]{l} \\catchOnError\\ onErrorResumeNext\\ retry\\{\color{white} a}\end{tabular}}                                                                            \\ \hline
\multicolumn{1}{|l|}{\textbf{Utility}}                  & \multicolumn{1}{l|}{\begin{tabular}[c]{l}\\ observeOn\\ subscribeOn\\ publish\\ share\\ ofType\\ doOnNext\\ doOnError\\ doOnCompleted\\  toIterable\\ toList\\{\color{white} a}\end{tabular}}     \\ \hline
\end{tabular}
\end{table}


\subsection{Formalizing Operators}

So far we have seen which are the most useful operators and how they aid in making our reactive library more useful from a user perspective. Things get more complicated when we try to analyze them from a theoretical point of view: operators can be viewed as state machines, containing an internal state which is modified whenever an event occurs, following the semantics of the specific operator. 

For this reason, trying to formalize them as a functional and pure structure becomes very difficult, resulting in more confusion than clarity. This outcome should not come as a surprise: operators defined on other more common data structures such as lists or trees are state machines as well, usually hiding their state using function parameters:

\hminted
take :: Int -> [a] -> [a]
take _ []     = []
take 0 _      = []
take n (x:xs) = x : take (n-1) xs
\end{minted}

As we can see from this example implementation of take on lists, the internal state of the operator - i.e. the number of elements to be taken, \code{n} - is wired in the definition of the function, eliminating the need for an internal variable as a result.

This type of operator definition works very neatly for any pull-based data structure, where we can define the operators recursively on the structure of the collection. For \code{Observable} things become a little more complicated since we are dealing with a push-based collection, were elements are never gathered as a concrete collection in memory, not allowing, as a consequence, any type of structural recursion.

We will discuss in section \ref{future-work}, Future Work, how Event Calculus might be used as a technique to better define the semantics of operators for push-based collections.
\clearemptydoublepage

\listoftodos
\clearemptydoublepage

% And the appendix goes here
\appendix
%\import{chapters/appendices/}{app0A.tex}
\chapter{Appendix A}
\label{app:a}

\section{Duality - Category Theory}

Duality is an element of paramount importance in research we presented. Up until this moment we have only dealt with an informal definition of duality but its roots go very deep in the field of mathematics and category theory in particular.

Given the formal definition of a category, composed of objects $$A,\; B,\; C,\; ...$$ and arrows $$f,\; g,\; h,\; ... $$ and operations $$ dom(f) \quad cod(f) \quad id(A) \quad g \circ f $$ 

Given any sentence $ \Sigma $, we can create its dual $ \Sigma^{op} $ by interchanging $$dom \to cod$$ and the order of composition $$ g \circ f \to f \circ g $$ 

It is clear how $ \Sigma^{op} $ is a well formed sentence in the language of category theory $ CT $.

The principle of - formal - duality then tells us that for any sentence $ \Sigma $ in the language of category theory, if $ \Sigma $ follows from the axioms defined for categories, then so does $ \Sigma^{op}$. 

$$ CT \vdash \Sigma \implies CT \vdash \Sigma^{op} $$

From a visual perspective, this definition boils down to reversing the order of composition of the arrows. 

\begin{center}
\begin{minipage}{.3\textwidth}
\begin{tikzpicture}
\large
  \matrix (m) [matrix of math nodes,row sep=3em,column sep=4em,minimum width=2em]
  { A & B \\
       & C \\};
  \path[-stealth]
    (m-1-1) edge node [above] {$f$} (m-1-2)
                 edge node [below] {$g \circ f\qquad$} (m-2-2)
    (m-1-2) edge node [right] {$g$} (m-2-2);
\end{tikzpicture}
\end{minipage}
\begin{minipage}{.3\textwidth}
\begin{tikzpicture}
\large
  \matrix (m) [matrix of math nodes,row sep=3em,column sep=4em,minimum width=2em]
  { A & B \\
       & C \\};
  \path[-stealth]
    (m-1-2) edge node [above] {$f$} (m-1-1)
    (m-2-2) edge node [below] {$f \circ g\qquad$} (m-1-1)
                 edge node [right] {$g$} (m-1-2);
\end{tikzpicture}
\end{minipage}
\end{center}

\section{Products \& Coproducts}

In category theory, the product of two objects in a category is the most general object which admits a morphism to each of the ones that compose it. The notion of product aims at capturing the essence of more specialized definitions found in various categories and areas of mathematics. The easiest way understand this construct is to start from the cartesian products in the category of sets.

Given sets $A$ and $B$, let us define the cartesian product as the set $$A \times B\; =\; \{\; (a, b)\; |\; a \in A,\; b \in B\; \} $$ There are two coordinate projections
\begin{center}
\begin{tikzpicture}
  \matrix (m) [matrix of math nodes,column sep=4em,minimum width=2em] { 
	  A & {A \times B} & B \\
	  & & \\
  };
  \path[-stealth]
    (m-1-2) edge node [above] {$fst$} (m-1-1)
                 edge node [above] {$snd$} (m-1-3);
\end{tikzpicture}
\end{center}
where $$ fst\; (a, b) = a \qquad snd\; (a, b) = b $$ 

It follows that given any element $ c \in A \times B $ $$ c\; =\; (\; fst\; (c),\; snd\; (c)\; )$$ 
The following diagram captures the essence of cartesian products.

\begin{center}
\begin{tikzpicture}
  \matrix (m) [matrix of math nodes,row sep=7em,column sep=6em,minimum width=2em] { 
	 {} & 1 & {} \\
	  A & {A \times B} & B \\
  };
  \path[-stealth]
  	(m-1-2) edge node [left] {$a\quad$} (m-2-1)
                 edge node [right] {$\quad b$} (m-2-3)
                 edge node [right] {$(a, b)$} (m-2-2)
    (m-2-2) edge node [below] {$fst$} (m-2-1)
                 edge node [below] {$snd$} (m-2-3);
\end{tikzpicture}
\end{center}

The definition of categorical products can be derived from generalizing the elements in the previous definition.

Coproducts are the dual notion of categorical products, representing the least general object to which the objects in the family admit a morphism. Within the context of set theory, the represent the disjoint union of sets.

Given sets $A$ and $B$, let us define their disjoint union as the set $$A + B\; =\; \{\; (a, 1)\; |\; a \in A\} \cup (b, 2)\; |\; b \in B\} $$ There are two injection functions
\begin{center}
\begin{tikzpicture}
  \matrix (m) [matrix of math nodes,column sep=4em,minimum width=2em] { 
	  A & {A + B} & B \\
	  & & \\
  };
  \path[-stealth]
    (m-1-1) edge node [above] {$left$} (m-1-2)
    (m-1-3) edge node [above] {$right$} (m-1-2);
\end{tikzpicture}
\end{center}
where $$ left\; (a) = (a,\; 1) \qquad right\; (b) = (b,\; 2) $$ 

The essence is captured by the following diagram, where we have simply reversed the arrows.

\begin{center}
\begin{tikzpicture}
  \matrix (m) [matrix of math nodes,row sep=7em,column sep=6em,minimum width=2em] { 
	 {} & 1 & {} \\
	  A & {A + B} & B \\
  };
  \path[-stealth]
  (m-2-1) edge node [left] {$f\quad$} (m-1-2)
  (m-2-3) edge node [right] {$\quad g$} (m-1-2)
  (m-2-2) edge node [right] {$[f,g]$} (m-1-2)
  (m-2-1) edge node [below] {$left$} (m-2-2)
  (m-2-3) edge node [below] {$right$} (m-2-2);
\end{tikzpicture}
\end{center}

where $$ [f, g](x,\delta) = \begin{cases} f(x), & \delta = 1 \\ g(x), & \delta = 2 \end{cases} $$

Withing the context of our research and the Haskell programming language, tuples represent products and the Either monad represents coproducts.

\section{Functors}

In category theory a functor is a mapping from one category to another or a homomorphism of categories where certain laws hold. 

Functors span thorough a large number of categories and are more and more common in modern programming languages. In Haskell the Functor class and its laws are defined as follows.

\hminted
class Functor f where
    fmap :: (a -> b) -> f a -> f b
  
fmap id      = id
fmap (p . q) = (fmap p) . (fmap q)
\end{minted}

\section{Covariance \& Contravariance}

\section{Continuations}

\section{Curry \& Uncurry}

Currying is the technique of transforming a function taking multiple arguments as input in one taking only the first argument and returning a function that takes the remainder of the arguments and returns the result of the initial function. Uncurrying is the opposite of currying, taking a curried function into one that accepts multiple arguments as input.

\hminted
curry   :: ((a, b) -> c) -> (a -> b -> c)
uncurry :: (a -> b -> c) -> (a, b) -> c

f :: a -> b -> c
g :: (a, b) -> c

f = curry g
g = uncurry f

-- f x y = g (x, y)
-- curry . uncurry = id
\end{minted}










\
\clearemptydoublepage

%
% file: localoperator.tex
% author: Victor Brena
% description: Briefly describes properties of the local operator.
%

\chapter{Appendix B}
\label{app:b}

\section{Formal Rx Implementation}

The following code presents a continuation based reference Reactive Library where each additional feature is implemented as discussed in the respective formalization section in Chapter \ref{chap:outoftherabbithole} .

\haskellcode[lastline=212]{src/rx-semantics-ReaderT-ContT-final.hs}
\clearemptydoublepage

% Apparently the guidelines don't say anything about citations or
% bibliography styles so I guess we can use anything.
\backmatter
\bibliographystyle{siam}
\refstepcounter{chapter}
\bibliography{thesisbiblio}
\clearemptydoublepage
%
% Add index
%\printindex
%   
\end{document}